\documentclass[10pt]{article}
\usepackage[margin=20mm]{geometry}
\usepackage{amsmath,amssymb,fontspec,unicode-math,xurl,hyperref}
\IfFontExistsTF{DejaVu Serif}{\setmainfont{DejaVu Serif}}{\setmainfont{Latin Modern Roman}}
\setmathfont{latinmodern-math.otf}
\setlength{\parindent}{0pt}\setlength{\parskip}{6pt}\sloppy
\begin{document}
ELVIA scientific verification report\par
0053 | I25-40159\par
Position: Certified Robustness Does Not (Yet) Imply Model Security\par
Outcome: Minor; two bounded technical-background inaccuracies; no demonstrated defect in the inspected core non-implication or norm conversion\par
Human scientific decision: Pending\par
\section*{1. Identity, scope and outcome}
The reviewed artifact is the 14-page publisher-linked ICML 2025 paper by Andrew C. Cullen, Paul Montague, Sarah Erfani and Benjamin I. P. Rubinstein [E1]. Its SHA-256 is 3f96194f25866608a43fa8e673f6e6e684ea024b36a0833e3e7f8aa22b04a6a4. The original bytes match the source-acquisition record. Printed pages 1-14 equal the 1-based PDF pages.\par
All main-text pages 1-9 were read. Two categorical descriptions of existing certification methods require bounded corrections. The core distinction between local prediction invariance and broader security, and the norm-conversion factor in Theorem 2, survived the checks below. No new experimental reproduction is claimed.\par
\section*{2. Supported findings and validated corrections}
F1 — Training is not a prerequisite for computing these certificates (Minor). Section 2.2, p3, states that exact approaches require substantial training changes; Section 2.2.2 then makes an augmented training loss a requirement for using IBP or convex-relaxation bounds. This conflates computing a sound bound on a fixed network with training a network to obtain useful bounds. The cited Xu et al. (2020), Section 3.1, treats a supplied computational graph and fixed unperturbed parameters; robust-loss training is a separate application in Section 3.3 [E3].\par
A complete admissible counterexample uses the fixed network h(t)=ReLU(t), with logits f(t)=(1-h(t),h(t)), input t=1/5 and radius 1/10. The input interval is [1/10,3/10], on which h=t. Ordinary interval propagation gives class-0 logits in [7/10,9/10] and class-1 logits in [1/10,3/10]. Their worst-case separation is 2/5>0, so class 0 is certified throughout that interval. No augmented training objective is required for this computation. Exact Fraction arithmetic independently confirms these bounds in analytic\_check\_results.json.\par
The strongest alternative reading is that specialized training is often needed for practically useful tightness, especially for deep-network IBP. That qualified statement is compatible with [E3], but differs from the paper’s categorical prerequisite. The verified correction is to separate post-training verification from training for larger or tighter certificates. The example does not establish useful radii, speed or scalability for arbitrary networks. Minor is supported because the false premise affects the background comparison of methods; it is not used in the proof of either central non-implication or norm conversion.\par
F2 — Non-ReLU certification predates the cited recent advance (Minor). Section 2.2.2, pp3-4, describes applicability beyond ReLU as only recently demonstrated, citing Shi et al. (2023). Yet its own Zhang et al. (2018) reference develops CROWN for general activation functions. Section 3.2 and Table 5 explicitly cover tanh, sigmoid and arctan, with nonzero certified bounds [E4]. Section 4 also applies CROWN to pretrained ReLU models.\par
The strongest alternative is that a particular architecture, tightness improvement or practical scale was intended. The passage does not specify such a restriction. A source-supported correction is to acknowledge earlier general-activation certification and identify the narrower advance meant by the 2023 citation. This repairs the historical and capability description; it does not show that all modern architectures are easy to certify. Minor follows from the limited dependency: the paper’s broad concerns about deployment costs and threat-model scope do not require a claim that non-ReLU certification began recently.\par
\section*{3. Core mathematical checks and rejected stronger readings}
Theorem 1, pp4-5: a valid radius is a prediction-invariance guarantee, not a semantic correctness or input-history guarantee. The constant-probability illustration can show certified misclassification, but its constant predicted class does not itself exhibit a class-flipping attack under Eq. (1). This incompleteness of the illustration does not refute the non-implication.\par
A complete independent example takes S=[0,1], F(t)=0 for t<1/2 and F(t)=1 otherwise, with semantic truth y(t)=0 throughout. The clean input 1/4 is correctly classified and an attacker can replace it by 3/4, preserving truth but changing the prediction. Each point has valid radius 1/5, below the boundary distance 1/4. The attack magnitude is 1/2, outside the original certified radius. Thus the example demonstrates a certificate at an already-attacked input, without violating any within-radius guarantee. The paper’s p4 expressly acknowledges the real guarantee against correctly bounded attacks and permits heuristic risk ranking. Accordingly, its no-information wording is read as no intrinsic attack-status guarantee, not universal statistical independence.\par
Theorem 2, p7: for p,q>=1, a radius r\_p in the p norm transfers to the q-norm radius shown below. If q<=p, the p norm is at most the q norm. If q>=p, the p norm is at most d\textasciicircum{}(1/p-1/q) times the q norm. These establish inclusion of the transferred ball; one-coordinate and equal-coordinate vectors respectively make the bounds sharp in ambient space.\par
Beyond this radius there can be points outside the known certified ball. This does not prove an actual adversarial example for every classifier, and a restricted permissible domain can remove geometric witness directions. The theorem says potential attacks, its proof concludes noncoverage, and nearby prose says attacks may exist. These qualifications support the coverage reading. A constant classifier would refute a stronger universal attack-existence claim, but does not establish a defect in this stated, qualified argument. Merely changing the norm also does not revoke the original certificate.\par
Eqs. (3)-(4), p3: the lower top-class and upper competing-class probability bounds, inverse-normal signs and sigma/2 factor are consistent with the randomized-smoothing guarantee [E6]. The text explicitly identifies the smoothed classifier as the certified object. Sampling confidence and abstention remain relevant; no Monte Carlo procedure was rerun.\par
\[R_q=r_p\min\{1,d^{1/q-1/p}\},\qquad B_q(x,R_q)\subseteq B_p(x,r_p).\]\par
\section*{4. Security evidence, version comparison and remaining limits}
The certificate-guided-attack concern is conditional on access. The paper states this on p5; the cited Cullen et al. (2024), Sections 3.1 and 4.1, assumes relevant model/certificate access and uses certificates to guide search beyond known certified regions [E5]. This supports a conditional information advantage, not a contradiction of the certified-radius guarantee. No claim of universally worse security or measured user complacency is established. Section 6 acknowledges that partial improvements can enhance security.\par
The inspected publisher material and paper do not identify a new code or data release requiring reproduction for this position paper. Cited prior-work experiments were not rerun. The deployment, user-expectation and governance proposals were evaluated as arguments and research questions; no legal compliance opinion or new experiment was invented.\par
The latest retrieved arXiv record is v2 dated 11 August 2025 [E2]. Checked pp3-5 and 7 retain the relevant descriptions and arguments. A main-text word comparison found only small text/encoding differences; it is a locator aid, not proof of formula-level version equivalence. No repair to F1/F2 was identified in the checked passages. OpenReview history was inaccessible.\par
\section*{5. Conclusion and repair/adoption tracking}
The supported overall impact is Minor. The two corrections concern the description and historical scope of existing certification methods; the inspected core non-implication and norm conversion remain valid.\par
Item | Judgment | Repair verification | Adoption / remaining work\par
Special training required for certification | Minor | Fixed-network interval example and primary methods verified | Textual correction proposed; no adoption\par
Non-ReLU extensions only recent | Minor | 2018 CROWN primary source verifies earlier scope | Correct historical wording; no adoption\par
Local robustness versus broader security | No Issue Found in inspected claim | Admissible example and norm factor checked | No new attack or deployment validation\par
External experiments and full bibliography | Not fully verified | Selected primary-reference checks only | Broader reproduction remains unperformed\par
Human scientific decision | Pending | AI inspection only | publication status and human judgment pending\par
\section*{Supporting checks}
[E1] Cullen et al. (2025). PMLR 267:81185-81198. https://proceedings.mlr.press/v267/cullen25a.html.\par
Cullen et al., arXiv:2506.13024v2, 11 August 2025. https://arxiv.org/abs/2506.13024v2.\par
Xu et al. (2020). Automatic Perturbation Analysis for Scalable Certified Robustness and Beyond. NeurIPS 33:1129-1141. https://proceedings.neurips.cc/paper/2020/hash/0cbc5671ae26f67871cb914d81ef8fc1-Abstract.html. Relevant Sections 2, 3.1 and 3.3.\par
Zhang et al. (2018). Efficient Neural Network Robustness Certification with General Activation Functions. NeurIPS 31. https://proceedings.neurips.cc/paper/2018/hash/d04863f100d59b3eb688a11f95b0ae60-Abstract.html. Relevant Sections 3.2 and 4, Table 5.\par
Cullen et al. (2024). Et Tu Certifications: Robustness Certificates Yield Better Adversarial Examples. PMLR 235:9745-9761. https://proceedings.mlr.press/v235/cullen24a.html.\par
Cohen, Rosenfeld and Kolter (2019). Certified Adversarial Robustness via Randomized Smoothing. PMLR 97:1310-1320. https://proceedings.mlr.press/v97/cohen19c.html.\par
\section*{Supporting checks}
This script uses only exact rational arithmetic and produces analytic\_check\_results.json.\par
The original PDF was extracted with installed Poppler pdftotext -layout.\par
No author code was executed; no new training, model inference or paid API call was made. Full validation of every external citation and empirical deployment claims remains unperformed. Both proposed textual corrections are analytically or source verified, but author adoption and human scientific decisions remain pending. The supported overall impact is Minor.\par
\end{document}
